\chapter{Discussion}
\label{ch:discussion}

\section{The foveation result}
\label{sec:disc-finding}

Gaze-contingent input foveation at human acuity costs nothing measurable. The difference at
$f_{c0} = 40$ is at most $0.0071$~bits per fixation, under $1\%$ of what the sharp control gains over the
centerbias, and a loss appears once the blur is pushed past human acuity (\subsecref{results-arms}). At
the coarse cutoffs the map-level diagnostics assign the loss to spreading rather than to mislocation, the
maps carrying more entropy at $f_{c0} = 20$ and $10$ while AUC moves by at most $0.002$
(\subsecref{results-diagnostics}).

At human acuity the foveation removed nothing the fit depended on measurably. At the coarser cutoffs it
did. The model draws on detail the coarser blur removes, by about $2\%$ of its gain over the centerbias at
$f_{c0} = 10$, and the per-fixation cost deepens with saccade amplitude (\figref{stratified-amplitude}); the measured cost is an
upper bound on what the discarded detail was worth, provided the frozen backbone --- fitted to sharp
photographs --- is no better on foveated input than on sharp. Nothing here tests that; unfreezing the
backbone would (\secref{disc-limits}).

The fixed-centre arms separate degrading the periphery from letting the sharp region follow the gaze.
With the same blur held at the image centre, the arm sits below the gaze-contingent one at every cutoff,
and the difference between the two, the gaze-contingency term, grows from $0.003 \pm 0.003$~bits at
$f_{c0} = 40$ to $0.013 \pm 0.004$ at $10$ (\subsecref{results-gaze-contingency}). Degrading the periphery is what costs, and letting the sharp
region follow the gaze gives part of that cost back. What produces that return is not established.

\dg{}'s authors put a related question to the read-out (\secref{position}) and found that \enquote{the
distance to the fovea plays only a minor role for guiding free-viewing gaze on average}
\citep{kummerer2022deepgaze3}. Those variants were fed the same
sharp image, so the test held the available information fixed. The result here is the input-level
counterpart. At human acuity, changing the information itself leaves performance where it was, within
what this design can measure.

\section{The per-image diagnostic}
\label{sec:disc-per-image}

\dg{}'s fit varies more from image to image than any arm differs from another, and that variation
follows how far the fifteen subjects spread (\secref{results-per-image}). \subsecref{stats-plan} takes every difference inside an image for that
reason, and the diagnostic characterises what the pairing removes. Part of that spread belongs to the
metric (\secref{results-per-image}), so a low per-image score is not on its own evidence that the model
failed there.

\section{Limitations}
\label{sec:disc-limits}

Several limitations bound what the comparison resolves and how far the result reaches.

\begin{itemize}
  \item \emph{The backbone is frozen.} The features every arm reads were learned on sharp photographs,
    and only the read-out adapts to the foveated input. Unfreezing part of the backbone, or training it
    on foveated input from the start, is untested (\secref{disc-future}).
  \item \emph{The evaluation is scored on human scanpaths.} Every history the model conditions on is a
    human scanpath and every scored location is a human fixation (\subsecref{metrics}), so the foveated
    model's sharp region always sits where a subject was looking, and the model never has to find a target
    with its own degraded periphery. What the comparison measures is the value of peripheral detail for
    predicting the next human fixation from a human history. Whether a foveated model would produce
    human-like scanpaths on its own is a separate question this design does not ask.
  \item \emph{Native resolution is a shift for the backbone.} \dg{} was released fitted on resized
    stimuli, and the arms here read each stimulus at its own size so that the sharp region stays a disc
    (\subsecref{gp-model}). Every arm carries that shift equally, so it is not a difference between
    them, and nothing here tests whether it interacts with the foveation.
  \item \emph{Faithful foveation is mild on lab-sized images.}
    Because acuity falloff is set in degrees and an MIT1003 stimulus spans $29^\circ$ along its long
    side, the foveation at the human setting leaves a third of all saccade targets untouched.
    Differences large enough to read off without ten-fold pairing appear only at coarser-than-human
    strengths, so the argument rests on the shape of the dose--response and on the fixed-centre contrast
    rather than on the size of any single point.
  \item \emph{A single dataset.} All results are on MIT1003 free-viewing by adults; other datasets and
    tasks (visual search, other populations) may behave differently.
  \item \emph{Fifteen subjects, thinly described.} MIT1003 provides fifteen subjects for every
    image, described in a single sentence as males and females aged between 18 and 35
    \citep{judd2009mit1003}. Each likelihood in \chref{results} is therefore measured against those
    fifteen people rather than against the expected value more subjects would estimate, and age is
    given only as that range, while first language and script direction are not recorded at all, so
    none of them can be controlled for.
  \item \emph{The falloff shape is fixed.} Every foveated arm inherits the same half-acuity eccentricity
    $e_2 = 2.3^\circ$ (\subsecref{gp-model}). Human-strength here means human-strength for
    one acuity profile, and the thesis does not test others.
  \item \emph{The mechanism models acuity falloff only.} Resolution is not the only peripheral limit. An
    object resolvable in isolation becomes hard to identify when others lie near it, \emph{crowding}
    \citep{strasburger2011peripheral}, so degrading peripheral resolution reproduces part of peripheral
    vision and not all of it.
  \item \emph{Fixation duration is not represented.} \dg{} takes fixation coordinates only
    \citep{kummerer2022deepgaze3}, so if foveating the input changed how long subjects dwelled at each
    location without changing which locations they chose, this evaluation would not register it.
\end{itemize}

\section{Future work}
\label{sec:disc-future}

The most direct extension is replication. Every number reported here comes from one dataset, so the
dose--response of \secref{results-foveation} may be a property of the mechanism or a property of MIT1003.
Running the same arms and the same $f_{c0}$ sweep on a second dataset of adult free-viewing gaze would
separate the two. The evaluator already accepts any model matching the \dg{} forward signature
(Appendix~\ref{app:eval}), so the work is a loader and a centerbias fitted to the new stimuli
(\subsecref{metrics}). The dataset must record its viewing geometry, because without pixels-per-degree
the two cannot be foveated at a matched strength in degrees (\subsecref{gp-model}), leaving display
scale confounded with foveation.

Unfreezing part of the backbone would separate the detail the foveation discards from the DenseNet's
unfamiliarity with blurred input, which the cost at the coarse cutoffs mixes (\secref{disc-limits}).
Compute kept it out of scope, and it would change what the comparison measures, since a difference
between arms would then include what each backbone learned from its own input.

The longer-range question is developmental, whether a model's \emph{learning trajectory} recapitulates
the development of human gaze, which changes measurably with age
\citep{helo2014maturation, acik2010developmental}. This thesis makes no claim that a read-out's training
trajectory resembles that development; the front-end makes the comparison possible, and testing it needs
a developmental eye-tracking dataset that records its viewing geometry, which the replication above also
requires.

Smaller extensions follow from constraints already stated. A transformer host would test whether the dose--response
survives a change of generator (\secref{position}). A duration channel would let the evaluation register
a change in dwell time (\secref{disc-limits}).
